\documentclass[11pt,a4paper]{article}
\usepackage[T1]{fontenc}
\usepackage[utf8]{inputenc}
\usepackage{lmodern}
\usepackage[margin=25mm,headheight=15pt]{geometry}
\usepackage{microtype}
\usepackage{amsmath,amssymb,amsthm,mathtools}
\usepackage{booktabs,longtable,array,enumitem}
\usepackage{xcolor,listings,fancyhdr}
\usepackage{xurl}
\usepackage[colorlinks=true,linkcolor=teal!60!black,citecolor=teal!60!black,urlcolor=teal!60!black,bookmarksnumbered=true]{hyperref}
\hypersetup{pdftitle={Answering the Expert Questions for Leant2},pdfauthor={Research report prepared for Vladimir Reshetnikov},pdfsubject={Dependent program synthesis, carrier invention, conservative evaluation, and Lean-native search}}
\setlength{\parindent}{0pt}
\setlength{\parskip}{5pt plus 1pt}
\setlength{\emergencystretch}{2em}
\setcounter{tocdepth}{2}
\setlist{nosep,leftmargin=1.6em,topsep=4pt}
\renewcommand{\arraystretch}{1.16}
\pagestyle{fancy}
\fancyhf{}
\fancyhead[L]{\small\textsc{Leant2: Expert Questions}}
\fancyhead[R]{\small Research and implementation design}
\fancyfoot[C]{\small\thepage}
\renewcommand{\headrulewidth}{0.3pt}
\lstset{basicstyle=\ttfamily\footnotesize,columns=fullflexible,keepspaces=true,breaklines=true,frame=single,framerule=0.3pt,rulecolor=\color{black!25},backgroundcolor=\color{black!2},xleftmargin=4pt,xrightmargin=4pt,aboveskip=8pt,belowskip=8pt,showstringspaces=false}
\newcommand{\code}[1]{\texttt{\small\detokenize{#1}}}
\newcommand{\Nat}{\mathbb{N}}
\newcommand{\Bool}{\mathsf{Bool}}
\newcommand{\List}{\mathsf{List}}
\newcommand{\Fin}{\mathsf{Fin}}
\newcommand{\Type}{\mathsf{Type}}
\newcommand{\Prop}{\mathsf{Prop}}
\newcommand{\Comp}{\mathsf{Comp}}
\newcommand{\supp}{\mathsf{supp}}
\newcommand{\len}{\mathsf{length}}
\newcommand{\rev}{\mathsf{reverse}}
\newcommand{\foldr}{\mathsf{foldr}}
\newcommand{\nil}{[\,]}
\newcommand{\cons}{\mathbin{::}}
\newcommand{\app}{\mathbin{+\!+}}
\newcommand{\question}[2]{\subsection{#1. #2}\label{q:#1}}
\newcommand{\answer}{\textbf{Proposed answer.}\ }
\newcommand{\status}[1]{\textbf{Evidence status:} #1\par}
\newtheorem{proposition}{Proposition}[section]
\newtheorem{lemma}[proposition]{Lemma}
\newtheorem{definition}[proposition]{Definition}
\newtheorem{theorem}[proposition]{Theorem}
\theoremstyle{remark}
\newtheorem{remark}[proposition]{Remark}
\newcolumntype{L}[1]{>{\raggedright\arraybackslash}p{#1}}
\begin{document}
\begin{titlepage}
\vspace*{16mm}
{\color{teal!60!black}\large\textsc{Research report \quad / \quad Lean-native synthesis}}\par
\vspace{12mm}
{\Huge\bfseries Answering the Expert\\[4pt] Questions for Leant2\par}
\vspace{8mm}
{\Large Sound search, carrier invention, conservative evaluation,\\[3pt]
and a concrete implementation plan\par}
\vspace{14mm}
{\large Prepared for Vladimir Reshetnikov\par}
{\large September 21, 2026\par}
\vspace{12mm}
\hrule
\vspace{5mm}
\textbf{Scope.} An answer to all 31 questions in the nine research areas of
\emph{Leant2: Design for Further Improvements}, together with corrections to
several premises of those questions. The investigation uses the repository,
Lean 4.34.0 source and documentation, and primary research literature.
\vspace{5mm}

\textbf{Repository snapshot.}\par
\path{VladimirReshetnikov/Leant2}\par
\path{784664ff34e819422510a4d470ab07fe48bb736b}
\vspace{5mm}

\textbf{What is established here.} Mathematical arguments, source-verified
integration points, and a reproducible finite-state experiment. Proposed
algorithms are distinguished from implemented or measured Leant2 features.
The repository's benchmark results were not rerun, and the Lean examples in
this report were not compiled.
\vfill
{\small Main source: \path{docs/proposals/11-further-improvements/}.\\
The question identifiers R1.1--R9.2 retain the ordering of the source document.}
\end{titlepage}

\begin{abstract}
Leant2's further-improvements proposal identifies the right obstacles to a
practical dependently typed synthesizer, but some of its questions can now be
answered more concretely, and several proposed shortcuts require correction.
We investigate all 31 expert questions. Important existing ingredients include
Canonical-min's assignment-triggered constraint suspension, Smyth's
higher-order example support, AutoLifter's synthesis of auxiliary summaries,
Hoogle+'s dictionary-based treatment of type classes, and Lean's reusable
induction and library-suggestion interfaces. None is a drop-in solution to
Leant2's complete problem.

We propose a native-Lean constraint orchestrator with persistent branch
states, dependency-aware residual evaluation, bounded and backtrackable
summary invention, typed recursive-call capabilities, and staged proof
obligations. We separate five guarantees that the original plan sometimes
conflates: acceptance soundness, pruning soundness, bounded search
completeness, ranking optimality, and replay determinism. We prove a safe
finite-sample carrier lower-bound criterion, give a minimal four-state
example, explain why reverse is already a fold into lists, and show why
finite examples cannot generally determine a minimal useful carrier type.
The final sections give source-level integration points, an implementation
sequence, and an evaluation protocol. Model-guided suggestions remain
optional, typed, auditable proposals rather than part of the trusted core.
\end{abstract}

{\small\setlength{\parskip}{0pt}
\tableofcontents}
\clearpage

\section{Executive conclusions and scope}
\label{sec:scope}

The most useful next step is not to replace Leant2 with a completely new
synthesizer. It is to make the existing native-Lean search state explicit,
separate its guarantees, and introduce a small number of shared mechanisms:
a constraint-dependency service, typed partial observations, and a bounded
schema language for carriers, motives, and recursive programs. The difficult
invention problems should use those mechanisms rather than each growing a
separate search engine.

The source proposal is a design document dated September 18, 2026. Its own
measurements refer to commit \code{33cec8b}, Lean 4.34.0, and one Windows
workstation. This report instead fixes the inspected repository at commit
\code{784664f}; its conclusions about implementation refer to that snapshot,
not to hypothetical later changes. Reported baseline successes in the source
are evidence about that source's experiments, not new measurements made
here.\cite{leantrepo,leantproposal}

\begin{longtable}{@{}L{10mm}L{56mm}L{84mm}@{}}
\toprule
Area & Main answer & First implementation consequence\\
\midrule
\endhead
R1 & Shared constraints need persistent branch-local state, not a sum of independent goal costs. & Build explicit states and dependency tracking; begin with a safe zero lower bound.\\
R2 & Minimal finite machines are computable; minimal useful Lean carrier types are a different problem. & Use collision-directed, bounded summary synthesis with backtracking.\\
R3 & Higher-order live evaluation already exists; full dependent Lean support does not follow automatically. & Improve kernel-backed per-observation evaluation before adding a certified fragment evaluator.\\
R4 & Useful finite motive classes can be enumerated; unrestricted motive invention cannot be reduced to one finite list. & Reuse Lean's induction APIs and add targeted, dependency-closed generalization.\\
R5 & A recursive interpreter generally agrees with well-founded realizations propositionally, not by conversion alone. & Preserve equation certificates and prove interpreter--realizer correspondence.\\
R6 & Programs and proofs should exchange constraints, with expensive proving scheduled separately. & Admit local-context proof obligations and exploit recursion-aligned verification conditions.\\
R7 & Hoogle+ already handles type classes through dictionaries. & Extend a typed AND-hypergraph abstraction; use Lean's actual suggestion interface.\\
R8 & Fixed description cost is defensible, but it is not a theorem about user intent. & Separate immutable ranking from adaptive search; label observational buckets honestly.\\
R9 & No checked accuracy number answers Leant2's exact carrier-invention question. & Run a controlled study; make replay consume immutable typed proposals rather than rerun a model.\\
\bottomrule
\end{longtable}

\subsection{How to read the evidence}

Statements about existing systems below cite their papers or the relevant
versioned source. Algorithms introduced with ``propose'' or ``recommend'' are
designs for Leant2, not claims that an implementation already exists.
Mathematical propositions state their assumptions explicitly. A small Python
experiment checks finite instances of several arguments; it does not test
Lean or establish Leant2 performance. No outside experts were contacted as
part of this investigation.

Some literature changes the questions themselves. Smyth's implementation
supports higher-order examples, even though the presentation of its core
calculus initially restricts examples. Hoogle+'s implementation explicitly
translates type-class constraints into dictionary types. AutoLifter studies
precisely the invention of additional summaries needed to make a
compositional algorithm work. Canonical-min, published in March 2026,
describes suspended constraints that wake when a blocking metavariable is
assigned. These are substantive starting points, not merely neighboring
research areas.\cite{smyth,tygar,autolifter,canonicalmin}

\subsection{Corrections that affect the design}

Several corrections deserve emphasis before the detailed answers.
A function may be a fold into its result type yet require an expensive step;
reverse is an example. A carrier with a final decoding function is more
expressive than a direct fold whose state is the required output. A finite
set of values observed or generated at one depth is not an upper bound on a
Lean type's cardinality. Summing lower bounds for coupled goals can
 overestimate the remaining search cost. A final kernel gate cannot recover
solutions removed by an incorrect partial-program refuter. Finite agreement
on examples is not generally contextual equivalence, and a fixed wall-clock
deadline cannot universally provide machine-independent useful results.
These are mathematical distinctions, not implementation preferences.

\section{The five guarantees and a common semantic interface}
\label{sec:guarantees}

Let a query be a local telescope $\Gamma$, target type $T$, contract $C$, a
fixed environment $E$, and an axiom profile $A$. Let $\mathcal G_b$ be the
explicitly declared search grammar under structural bound $b$. A partial
state $s$ contains a term with holes, their well-typed contexts, universe and
term constraints, and any residual observations. Write $\Comp(s)$ for its
well-typed completions in the stated grammar.

\paragraph{Acceptance soundness.}
An accepted result contains $p:T$ and, when a contract was requested, a proof
of $C(p)$; both are kernel-checked in $E$ and their transitive axiom
dependencies satisfy $A$. This is the role of the existing acceptance gate,
which synchronously adds checked declarations in a scratch environment and
audits their axioms.\cite{leantgate}

\paragraph{Pruning soundness.}
A discarded partial state must satisfy
\begin{equation}\label{eq:prune}
 \forall p\in\Comp(s),\quad \neg C(p).
\end{equation}
More generally, a discarded branch can be inconsistent with its typing
constraints. This property is independent of acceptance soundness. A bad
refuter can make a completely sound accepting synthesizer return no answers.

\paragraph{Bounded completeness.}
There are two distinct bounded claims. \emph{Enumeration completeness} means
that every candidate in a finite declared grammar is eventually considered.
\emph{Certification completeness} additionally requires the promised
verification procedure to decide the relevant contracts or proof obligations.
A tactic timeout does not establish that a candidate is wrong, nor that no
proof exists. For general Lean propositions the report recommends explicit
``unproved'' outcomes rather than pretending that the second guarantee
follows from the first.

\paragraph{Ranking optimality.}
A returned term is optimal only relative to a fixed objective over a precisely
specified candidate set, with certified lower bounds for everything still
pending. Changing a learned search priority is permissible, but it does not
preserve an A* optimality argument for a different, immutable cost function
without additional bookkeeping.

\paragraph{Replay determinism.}
The same inputs, environment, proposal log, and logical work schedule should
reproduce the same event sequence and checked outputs. This is not the same
as finishing the same number of events before a wall-clock deadline.

A suitable internal result algebra therefore distinguishes a verified
candidate, a kernel-checkable refutation, exhaustion of a declared decidable
finite search, an unrefuted but unproved candidate, an interrupted search, and
an unsupported feature. It should never collapse the last three into
``uninhabited.'' The default interface can remain simple; the distinctions
belong in receipts and diagnostics rather than a proliferation of user
settings.

\subsection{A shared constraint architecture}

A proposed node is conceptually
\[
 s=(E,\Gamma,\Sigma,p,\mathcal H,\mathcal O,g,\rho),
\]
where $\Sigma$ is a persistent metavariable/constraint state, $\mathcal H$ is
the collection of typed holes, $\mathcal O$ is the collection of residual
observations and proof obligations, $g$ is immutable accumulated cost, and
$\rho$ is a stable replay identity. Scheduling scores are separate metadata.
A single dependency service answers which assignments may change which
constraints, observations, or tactic preconditions.

The first implementation should preserve the current depth-first order as
one queue policy. That isolates representation bugs from policy changes.
Persistent data structures reduce copying costs, but do not make arbitrary
cross-branch mutation safe. The inspected search already uses
\code{Meta.SavedState} around alternatives, while some search bookkeeping is
held in references outside that saved state. A first-class node must explicitly
classify every such reference as branch-local, immutable shared data, or
monotone diagnostic data.\cite{leantcore}

\section{R1: Search control with coupled goals}
\label{sec:r1}

\question{R1.1}{Can independent goals be exploited without copying everything?}

\answer Yes, by sharing immutable structure while representing assignments
and dependent continuations explicitly. However, independence must be tested
on the full constraint problem, not just on the visible goal expressions.
Aesop and Canonical provide important proof-search precedents; the newer
Canonical-min makes the assignment-triggered suspension mechanism especially
concrete. Its small dependent-type core is not a replacement for Lean's full
metaprogramming state.\cite{aesop,canonical,canonicalmin}

Define the support of a goal to include reachable unassigned term and universe
metavariables in its type, local declarations and their values, pending
constraints, and delayed assignments. Compute this transitively. For example,
if one goal mentions $\mathord?u$ and another $\mathord?v$, but their
declarations both depend on $\mathord?\alpha$, their visible metavariable
sets can be disjoint while the goals remain coupled.

There is a second coupling relation: a contract may relate otherwise
independent values. Holes $x:\Nat$ and $y:\Nat$ are typing-independent, but
the observation $x+y=7$ couples their acceptable assignments. The dependency
graph should therefore have constraint and observation nodes, with hyperedges
to all relevant holes. Shared rigid parameters do not by themselves force
joint search; shared writable constraints do.

A conservative factorization criterion is disjoint writable support after
conditioning on a frozen shared interface. Under that criterion, assignments
from the factors commute. Treat an opaque tactic that may inspect or assign
arbitrary reachable metavariables as touching their entire support unless it
has a narrower declared effect. This sacrifices some parallelism rather than
risking cross-branch interference.

For reuse, cache a checked solution \emph{schema} abstracted over a canonical
local telescope. Instantiate it freshly and recheck it in the destination
state. Do not transplant raw local or metavariable identifiers. For data
holes, retain a lazy stream of solutions: reusing only the first inhabitant
can lose every solution of the surrounding behavioral contract. A timeout
or failed tactic attempt is not a reusable negative theorem. Proven
inconsistency or exhaustion of an identical finite constraint problem is.

\textbf{Concrete recommendation.} Add a persistent dependency graph and wakeup
queue around native Lean metavariables. Keep branch snapshots for mutable
assignments, share syntax and closed lemmas, and factor only components that
pass the conservative test. This is a principled middle ground between
whole-tree duplication and a permanently fixed sequential order. It does not
eliminate the intrinsic combinatorics of genuinely coupled choices.

\question{R1.2}{What admissible heuristic should be used?}

\answer Start with $h(s)=0$. Introduce stronger lower bounds only with an
explicit cost-preserving abstraction argument.

The proposed ``one unit per open non-leaf goal'' is not generally admissible.
An exact local can have zero assigned production cost. Moreover, one
unification-producing refinement can close several open constraints. If two
constraints each require the same assignment $\mathord?n:=0$, summing two
unit lower bounds counts the same future action twice. A lower bound on each
constraint separately is not automatically additive.

\begin{proposition}[Abstract-distance lower bound]\label{prop:heuristic}
Let the concrete search have nonnegative transition costs. Suppose an
abstraction maps every concrete solution state to an abstract solution and
every concrete transition to an abstract path whose cost is no greater.
Then the minimum abstract solution cost from the image of a state is a lower
bound on its concrete remaining solution cost.
\end{proposition}
\begin{proof}
Every concrete solution path maps to an abstract solution path of no larger
cost. The minimum over all abstract paths is no larger than the cost of each
such mapped path, hence no larger than the minimum concrete solution cost.
\end{proof}

A practical abstraction forgets selected dependent indices and behavioral
constraints while retaining a finite typed production hypergraph. Forgetting
constraints must add possibilities, not remove them. Zero-cost unification
and exact closure must still be possible in the abstraction. For individually
valid lower bounds $h_i$, $\max_i h_i$ is safe. Their sum requires disjoint
work or a proved partition of transition costs among abstractions.

Finite zero-cost closure should be saturated before queue insertion, or the
search should impose a positive secondary structural bound. Otherwise
infinitely many zero-cost rewrites or casts can prevent fair progress even
when the primary cost is nonnegative. The same caution applies to heuristic
term-size measures that deliberately ignore proof arguments.

\question{R1.3}{Can just-in-time learning preserve soundness and stable ranking?}

\answer It can preserve acceptance soundness and eventual enumeration while
changing the finite-budget answer set. Stable ranking requires a separate,
fixed ranking model. Probe establishes a useful precedent for learning from
partial successes during enumeration, not a guarantee that user-visible
answers remain invariant under such adaptation.\cite{probe}

Use two scores. The search score can reward constructions that satisfy more
observations or reduce important blockers. The audit score is the fixed
objective used for displayed ranking and any optimality certificate. If
adaptive scores change, either rekey affected queue entries or explicitly
use them only for a separate guidance queue. Reserve a deterministic share
of expansions for the unguided grammar queue. This prevents starvation and
keeps every bounded alternative reachable, subject to the stated verifier
limitations.

Learning updates should occur at deterministic logical epochs rather than
elapsed-time thresholds. Record their inputs and resulting parameters. A
stable tie-breaker must depend on canonical identifiers, not hash-table
iteration order or fresh names derived from time. These choices make a run
replayable. They do not make it insensitive to stopping after different
numbers of steps.

If an adaptive rule introduces a new carrier type rather than reordering an
existing grammar production, it has enlarged the search space. That is
sometimes desirable, but it must be recorded as a grammar change rather than
justified by a ``reordering preserves completeness'' argument.

\question{R1.4}{Can a wall-clock deadline be machine-independent?}

\answer Not in general while also requiring useful progress to be returned.
Suppose discovering a solution requires $k$ sequential logical steps. A fast
machine completes them before the deadline and a slow machine does not. An
algorithm that returns the discovered solution on the fast machine cannot
always return the same result on the slow one. A constant empty answer would
be machine-independent but defeats the intended requirement.

Use deterministic logical work limits for reproducibility and a wall-clock
limit as an interrupt boundary. A receipt records the completed event prefix,
pending frontier, grammar version, and verified candidates. Deterministic
epoch boundaries can stabilize presentation, but the number of completed
epochs still depends on speed. Startup calibration is an approximation, not
a proof of equivalence: tactic, normalization, and I/O costs vary differently
across machines and queries.

Also, the universal Luby restart guarantee concerns a specified Las Vegas
restart setting. It does not directly prove optimality for correlated,
deterministic Leant2 lanes sharing examples and partial work.\cite{luby}
A deterministic portfolio with measured quotas is a reasonable engineering
choice; its performance should be measured rather than inferred from a
nonmatching restart theorem.

\section{R2: Carriers, accumulators, and auxiliary state}
\label{sec:r2}

This area benefits most from separating semantic questions from cost-bounded
program construction. Existing fold characterizations and auxiliary-summary
synthesis provide more guidance than the original proposal suggests, but
neither identifies a unique intended carrier from a few examples.

\question{R2.1}{Can finite examples determine a minimal carrier type?}

\answer There is an exact finite-state problem with a constructive solution,
and there is a substantially different Lean program-synthesis problem. A
useful implementation should expose the distinction internally.

For lists over $A$, write a fold realization as
\begin{equation}\label{eq:fold-realization}
 q(\nil)=z,\qquad q(a\cons xs)=s(a,q(xs)),\qquad h(xs)=d(q(xs)),
\end{equation}
where the carrier is $K$, $z:K$, $s:A\to K\to K$, and $d:K\to T$ is a
finishing or decoding function. A \emph{direct} fold into the result type has
$K=T$ and $d=\mathrm{id}$. These are not the same synthesis problem.

There are at least three optimization objectives: the number of reachable
semantic states, the syntactic cost of the carrier type, and the total cost
of the realized program, its proof, and its execution. They need not agree.
For example, $\Nat$ can encode arbitrarily large finite tuples, but encoding
and decoding them can defeat the purpose of introducing a simple product
carrier. An optimization over type syntax alone can choose a small type with
an enormously expensive step function.

\subsubsection*{The finite-state problem is decidable}

Suppose $A$ has $m$ elements, the allowed output set has $r$ elements, and
$K=\Fin(k)$ with $k\geq1$. There are exactly
\begin{equation}\label{eq:machine-count}
 k\,k^{mk}r^k=k^{mk+1}r^k
\end{equation}
choices of initial state, total transition table, and output table. Every
finite set of input--output observations can therefore be checked against
all $k$-state machines. Trying $k=1,2,\ldots$ finds the minimum number of
states consistent with the observations.

For a consistent finite sample this search does eventually succeed: build a
prefix trie of the \emph{reversed} sample words, add a sink state, and label
sample-ending states with their required outputs. Reading reversed words
matches the state-transition order of a right fold. Unspecified outputs can
be assigned any default. This proves existence of a finite sample-consistent
machine, not correctness on unobserved words.

The analogous bounded Lean problem becomes decidable only after specifying
an actually finite grammar of carrier types and implementations, finite
literal and universe choices where relevant, and a fragment whose typing
and observation checks are performed by total decision procedures. Merely
bounding the printed size of a carrier while permitting arbitrary helper
programs or unresolved proof searches is not such a finite decision problem.
A fair enumerator can instead report candidates and unresolved obligations
without claiming to decide them all.

\subsubsection*{A sound lower bound from finite observations}

For words $u,v$, a prefix $c$ is a distinguishing context when the sample
contains both $h(c\app u)$ and $h(c\app v)$ and their outputs differ. Define
the incompatibility graph on selected suffixes by adding an edge whenever
there is such a context.

\begin{proposition}[Observed-context lower bound]\label{prop:carrier-bound}
Every deterministic realization \eqref{eq:fold-realization} consistent with
the sample assigns different states to adjacent suffixes. Consequently the
number of reachable states is at least the chromatic number of the
incompatibility graph, and in particular at least its clique number.
\end{proposition}
\begin{proof}
For each prefix $c$ there is a function $S_c:K\to K$ with
$q(c\app u)=S_c(q(u))$, obtained by composing the fold transitions for the
letters of $c$. If $q(u)=q(v)$, then
$d(S_c(q(u)))=d(S_c(q(v)))$. This contradicts any observed distinguishing
context. The assignment $u\mapsto q(u)$ is therefore a proper coloring of
the graph by reachable states.
\end{proof}

This is a finite-observation consequence of the familiar residual-state view
of automata. Angluin's exact-learning result is relevant, but its query and
counterexample oracles provide information that a fixed handful of examples
does not.\cite{angluin}

Two implementation cautions are important. First, equality on the
\emph{known} entries of incomplete observation rows is not an equivalence
relation. The rows $(0,?)$, $(?,0)$, and $(1,0)$ are pairwise compatible along
the first two links but the first and third conflict. Blind union--find
merging is therefore unsound. Secondly, the number of colors found by a
greedy coloring algorithm is an \emph{upper} bound on chromatic number, not
a certified lower bound. A found clique is a safe lower-bound certificate.
State merging must also respect transition congruence; a coloring by itself
does not construct a valid fold machine.

\subsubsection*{A minimal four-state example}

Let $A=\{0,1\}$ and let $h(xs)$ mean that both symbols occur in $xs$.
Use suffixes $\nil,[0],[1],[0,1]$ and contexts $\nil,[0],[1]$. Their profiles
are as follows; entries give $h(c\app u)$.

\begin{center}
\begin{tabular}{@{}lccc@{}}
\toprule
Suffix $u$ & $c=\nil$ & $c=[0]$ & $c=[1]$\\
\midrule
$\nil$ & false & false & false\\
$[0]$ & false & false & true\\
$[1]$ & false & true & false\\
$[0,1]$ & true & true & true\\
\bottomrule
\end{tabular}
\end{center}

Every pair of rows differs in a column. Thus these observations force a
four-clique, and no carrier with fewer than four reachable states works.
There are nine distinct input lists among the twelve table entries. A
four-state carrier is $\Bool\times\Bool$, recording whether a zero and a
one have been seen. Start at $(\mathrm{false},\mathrm{false})$, set the
corresponding flag on each step, and finish by conjunction. Induction on the
list proves that each flag has precisely the stated meaning. This gives an
all-input implementation and proves that the lower bound is attained.

The accompanying experiment exhausts every one-, two-, and three-state
machine and finds no sample-consistent machine. Its counts are 2, 128, and
17,496, respectively, as predicted by \eqref{eq:machine-count}. It also
checks the four-state realization on 2,047 binary words of length at most
10. These tests corroborate the construction; Proposition~\ref{prop:carrier-bound}
and the flag invariant supply the proofs.

\subsubsection*{What the test cannot refute}

The proposal's suggestion to count the values of $K$ constructible from the
inputs at the current search depth is unsafe as a semantic cardinality
argument. A deeper step, a hidden encoding, or a different literal can reach
other values. Such a count can support a \emph{bounded grammar} refutation
only after proving that the enumeration exhausts every value relevant to
that grammar. It does not show that the type $K$ itself is too small.

Likewise, a direct-fold counterexample does not refute the same type when an
arbitrary finish function is allowed. Take $h=\mathsf{tail}$, with
$\mathsf{tail}(\nil)=\nil$. Then
\[
 h(\nil)=h([b])=\nil,\qquad h([a])=\nil\ne[b]=h([a,b]).
\]
A direct fold into lists would have to apply the same step to the two equal
states, so it cannot realize tail. Yet with the \emph{same} list carrier,
$q=\mathrm{id}$, $s=\mathsf{cons}$, and $d=\mathsf{tail}$ give a valid
factorization. A bounded, inexpensive finish grammar is thus part of a
meaningful carrier-invention specification.

\question{R2.2}{What replaces an alphabet for polymorphic functions?}

\answer For an explicitly relationally parametric list function, input
positions are sufficient; element equality is not automatically observable.
With an equality dictionary or additional operations, the model must retain
that extra structure.

Assume a pure parametric fragment in which
$f_A:\List(A)\to\List(A)$ satisfies naturality:
\begin{equation}\label{eq:naturality}
 f_B(\mathsf{map}(g,xs))=\mathsf{map}(g,f_A(xs))
 \quad\text{for every }g:A\to B.
\end{equation}
This is the relevant free-theorem assumption.\cite{wadler}
Let $e_n$ be the list of all positions in $\Fin(n)$ in order. For any list
$xs:\List(A)$ of length $n$, there is a function $g:\Fin(n)\to A$ mapping
positions to their entries, even when entries repeat. Naturality yields
\[
 f_A(xs)=\mathsf{map}(g,f_{\Fin(n)}(e_n)).
\]
Thus $f_{\Fin(n)}(e_n)$ is a word of selected positions describing the result
on every length-$n$ input. Using $\Fin(0)$ handles the empty type without
assuming an arbitrary function $\Nat\to A$ exists. In a fixed higher
universe, use the corresponding universe lift of this finite position type.

\begin{proposition}[Length-bounded canonical testing]
If two functions both satisfy \eqref{eq:naturality} and agree on $e_n$ for
each $0\leq n\leq N$, then they agree on every input of length at most $N$ at
every type covered by that naturality assumption.
\end{proposition}
\begin{proof}
Represent an arbitrary input as $\mathsf{map}(g,e_n)$ and apply naturality to
both functions. Their equal position words have equal images under $g$.
\end{proof}

This is input-length-bounded equality, not equality at every length. It also
does not follow merely from seeing a Lean type quantified over \code{Type}.
For instance, classical decidability permits the noncomputable family that
returns $xs$ when $A$ is subsingleton and returns $\nil$ otherwise. Mapping a
singleton list from a one-element type into a two-element type violates
naturality. A syntax and axiom-profile discipline must establish the
parametricity premise before exact canonical-test claims are made.

An explicit equality operation permits dependence on equality patterns of
positions. Equality-sensitive data-word or register-style models are then
more appropriate than a finite alphabet of concrete test values. A comparison
or algebraic dictionary carries still more information, and relational
reasoning must preserve its operations and assumed laws. In Leant2, use
canonical positions as a high-value test heuristic immediately, but permit
semantic quotienting only for a fragment with the required theorem.

\question{R2.3}{Can fusion be run backward to invent auxiliaries?}

\answer Yes as a constrained synthesis strategy, but not as a generally
complete inverse operation on a few examples. AutoLifter is directly relevant:
it synthesizes auxiliary information and combination functions for
lifting a reference computation into divide-and-conquer-like schemes. Its
minimality property is not global minimization over Lean carrier types, and
its discussion explicitly identifies incompleteness caused by greedy
choices. Its setting also supplies more structure than arbitrary few-example
Lean queries.\cite{autolifter}

The proposed Leant2 variant is \emph{collision-directed summary lifting}.
Maintain candidate summary features
\[
 q(xs)=(\phi_1(xs),\ldots,\phi_j(xs)).
\]
When observed contexts require two suffixes with equal summaries to be
separated, synthesize an additional feature from a bounded typed grammar
that distinguishes them. A finite table of candidate features turns feature
selection into a weighted covering problem over these collisions. This is
proposal generation, not a semantic proof that untested pairs are separated.

Next synthesize a transition that updates all features, and a finishing
function. If the transition needs information not retained by the summary,
that failure generates another lifting obligation. Crucially, retain
alternative feature sets and backtrack: a cheap first summary can be
impossible to update within the grammar even though another summary works.
The full candidate is certified by the invariant
\[
 q(\nil)=z,\qquad q(a\cons xs)=s(a,q(xs)),
 \qquad C(d\circ q),
\]
or, when a reference implementation $h$ is available, by the stronger
$d(q(xs))=h(xs)$ theorem. Learned collisions and feature-selection scores
never replace these obligations.

\subsubsection*{Reverse: a necessary correction and a useful cost example}

The fold characterization literature explicitly notes that reverse is a
fold.\cite{whenfold} Indeed,
\begin{equation}\label{eq:reverse-fold}
 \rev(xs)=\foldr\bigl((a,r)\mapsto r\app[a],\nil,xs\bigr).
\end{equation}
The base case is immediate, and the induction step is exactly
$\rev(a\cons xs)=\rev(xs)\app[a]$. So no semantic carrier lower bound can
show that reverse needs a carrier other than $\List(A)$.

What fails is an \emph{efficiency or syntax restriction}. Under the usual
linked-list cost model, appending to a growing reversed suffix costs
$0+1+\cdots+(n-1)=\Theta(n^2)$. A continuation carrier
$K=\List(A)\to\List(A)$ supports the fold
\[
 z=\mathrm{id},\qquad s(a,k)=\lambda acc.\,k(a\cons acc),
 \qquad d(k)=k(\nil).
\]
Its invariant is
\begin{equation}\label{eq:reverse-cont}
 q(xs)(acc)=\rev(xs)\app acc.
\end{equation}
For $a\cons xs$, the left side becomes $q(xs)(a\cons acc)$; the induction
hypothesis and associativity of append give
$\rev(xs)\app[a]\app acc=\rev(a\cons xs)\app acc$.
The fold builds one closure per element and its final application performs
one cons per element, giving linear work under that explicit cost model.
This illustrates why carrier invention should optimize the whole schema,
not seek a semantic explanation for every shallow-search failure.

For tree flattening, a similar continuation or difference-list summary can
avoid repeated copying. For Fibonacci, a pair can retain adjacent values.
Such templates belong in a shared bounded schema grammar; their usefulness
is demonstrated by invariants and cost models rather than by declaring the
unlifted type impossible.

\question{R2.4}{Is there a complete manageable motive grammar for primitive recursion?}

\answer There is a finite, complete enumeration \emph{relative to an explicitly
finite template class}. There is no reason to expect one small fixed list to
cover arbitrary parameterized primitive-recursive definitions economically.

Begin from the input telescope and the desired result. For each selected
recursor, enumerate constant result motives, generalization of eligible
parameters into function space, bounded products or dependent pairs of
summaries, and bounded abstractions of index expressions. Reversion must
include dependencies so that the generalized telescope remains well typed.
For an accumulator $a:S$, the motive acquires a factor $S\to T$; an
index-dependent accumulator may instead need a dependent function type.
Tupled results and proof-carrying summaries can require $\Sigma$-types rather
than ordinary products.

A breadth- or cost-bounded traversal of this finite template grammar is
complete for the grammar by construction, provided checking and the bounded
subsearches have their stated decision guarantees. It is not complete for
all useful auxiliaries whose types can be written in Lean. A fair schedule
of expanding bounds offers a relative semidecision procedure, not an
interactive completeness guarantee. Measure how much each template enlarges
branching and how many new benchmark families it unlocks.

\question{R2.5}{Can proof-generalization heuristics choose carriers?}

\answer They transfer as evidence-generating heuristics, not as a complete
algorithm. Both problems ask which information must remain available when a
recursive or inductive step changes its context.

A failed attempt to use an induction hypothesis can identify an accumulator
whose value changes across the call. Generalize that parameter and its
necessary dependencies. A failed attempt to compute a next output from the
current summary can identify lost information; propose a product component.
A repeatedly unresolved equation can suggest an invariant or helper
specification. Preserve the actual failed constraint and observed environments
so that the proposal is more precise than ``try a larger type.''

In the Leant2 architecture, carrier selection, motive generalization, and
invariant invention should therefore call one schema service with different
views of the same failure. Bound the number and complexity of proposed
features; retain an unguided fallback; prove every proposed invariant or
helper contract before using it as a permanent pruning fact. An unproved
lemma can order exploration, but cannot become an assumption in an accepted
proof or justify deleting other solutions.

\section{R3: Evaluation of programs with holes}
\label{sec:r3}

\question{R3.1}{Does live bidirectional evaluation already handle higher-order examples?}

\answer Yes. Smyth's implemented system supports higher-order examples and
polymorphism, and its purpose includes synthesis without trace-complete
examples. The source proposal's blanket first-order description should be
corrected. Hazel's typed hole closures provide part of the operational
foundation. Neither fact supplies an already-verified implementation for
Lean's entire dependent calculus.\cite{smyth,hazel}

For Leant2, represent an encountered hole as a typed closure
\[
 H(m,\sigma):\tau_m[\sigma],
\]
where $m$ was declared in telescope $\Gamma_m$ and $\sigma$ is a well-typed
substitution from that telescope into the evaluation environment. The same
metavariable can be encountered with different values of dependent indices;
each closure is typed after its own substitution. Nevertheless all closures
must be explained by one assignment to $m$, not independently chosen outputs.

For a function-valued hole, finite application observations are the natural
interface: ``at arguments $v_1,\ldots,v_j$, this closure must produce these
results or satisfy these relations.'' This avoids requiring a complete
extensional representation of a function. It does not license identifying
functions merely because their observed applications agree. Repeated calls
on equal arguments must share constraints, and dependent results must be
compared at compatible types, with explicit transport when necessary.

The first difficulty in moving to dependent types is therefore not simply
that evaluation can encounter a function. It is maintaining typed contextual
substitution, index-dependent result constraints, and the correlation of all
occurrences of one unknown program. Implementing these on native Lean
expressions avoids a second approximation of the user's type theory.

\question{R3.2}{What is a suitable incremental, backtracking normalizer?}

\answer Reuse the assignment-triggered suspension pattern, but design a
branch-aware cache around the actual evaluation contract. Canonical-min is
a recent concrete example of dependent constraints suspended on a blocking
metavariable and resumed on assignment, with backtracking state management.
It is not evidence that a general, certified, hole-assignment cache for
Lean 4.34's kernel is available off the shelf.\cite{canonicalmin}

The immediate improvement is incremental \emph{observation} evaluation, not
a complete replacement normalizer. Each observation records its residual,
read dependencies, and status: established true, established false, or
unknown. Assigning a hole wakes only affected observations. The existing
code instead aggregates blockers across the residual check; this can be
refined without changing the evaluator of record.\cite{leantcore}

A cache key must distinguish the expression, local telescope, relevant
assignment state, environment, transparency regime, and instance choices.
Closed certified results can be shared broadly. Open results should carry
branch versions or validated dependency fingerprints. A raw metavariable
identifier is insufficient: after backtracking, an identifier or a related
name can refer to a different assignment history. In one branch $m=0$ and in
another $m=1$; a cached evaluation of $m=0$ must not survive solely because
its surface expression is unchanged.

A safe design keeps immutable cache entries with read sets and uses a
persistent branch token for validity. Reuse an ancestor entry only when all
its recorded dependencies still have the same relevant values. More
aggressive dependency minimization is an optimization to verify later.
Delayed assignments and universe changes belong in the dependency model,
not just ordinary expression-metavariable assignments. Failed evaluation,
unsupported reduction, and exhausted fuel remain unknown.

\question{R3.3}{What is the cheapest provably conservative evaluator?}

\answer First retain kernel-backed reduction and make its observations
smaller. Then add a small typed, proof-producing fragment evaluator rather
than reimplementing the complete kernel.

At query preparation time, decompose recognized conjunctions, finite
\code{List.all} checks, Boolean conjunctions required to be true, and
\code{decide} wrappers using proved equivalences. Preserve arbitrary
unrecognized propositions as opaque observations. A custom \code{BEq}
instance need not coincide with propositional equality; rewrite \code{x == y}
into $x=y$ only with the appropriate correctness theorem. For a disjunction,
retain alternatives: a failed disjunct does not refute the whole observation.

For an initial reflected fragment, choose Boolean and natural-number
operations, constructors, projections, and a small collection of list
recursors. Give the syntax a denotation into Lean and prove each evaluation
rule correct. A reduction result can carry a proof relating the denotation
to its residual or value. Unsupported constructs are delegated to Lean or
left unknown. Differential tests against kernel reduction are useful for
finding bugs, but do not prove a partial evaluator conservative.

The phrase ``under-approximate evaluation'' requires care. Under-approximating
\emph{established reductions} is safe: some true facts remain unknown.
Under-approximating the \emph{set of possible outputs} is generally unsafe
for refutation: the omitted output might belong to a valid completion.
Instead maintain an over-approximation of possible outcomes, or produce a
certificate that a false observation is stable under every admissible
completion.

\begin{proposition}[Certified stable residual refutation]\label{prop:residual}
Suppose a well-typed partial term $p$ has a reduction certificate, valid
under every admissible assignment to its holes, that an observation
$O(p)$ equals false. Then no completion satisfying that observation can
extend the state.
\end{proposition}
\begin{proof}
Instantiate the certificate with any admissible completion. Substitution
preserves the certified equality, so the completed observation is false.
This contradicts its required truth.
\end{proof}

This argument assumes that the representation and substitution are well
typed. In the current source, \code{instantiatePartial} explicitly permits
intermediate expressions whose metavariable local contexts no longer match,
for reduction only. That is not a demonstrated bug, but the transformation
must be included in the pruning-soundness argument; invoking a correct
kernel primitive does not by itself establish that an arbitrary preprocessing
transformation preserves all completions.\cite{leantcore}

Lean4Lean is relevant to specifying and checking Lean, but the latest examined
revision describes correctness proofs for parts of its typechecker and
metatheory, not a ready-made proof of a new Lean 4.34 partial normalizer.
Extracting a small evaluator from a separately stated typed fragment is a
more contained first verification project.\cite{lean4lean}

\question{R3.4}{Can top-down angelic evaluation terminate and remain complete?}

\answer Yes under explicit finite-domain and finite-grammar hypotheses;
otherwise only conditional or relative guarantees are justified. Burst
provides the important angelic-execution/refinement precedent, but its
algorithmic guarantees do not transfer simply by calling another evaluator
``angelic.''\cite{burst}

An angelic call stands for a set or relation of possible results. A branch
may be refuted only if \emph{no} admissible assignment of results can satisfy
the observations, using a sound over-approximation of what actual recursive
calls can return. Guessing one convenient result and later finding it
inconsistent rejects that guess, not every possible completion of the
program. Results for repeated identical typed calls must be correlated.

For a finite program grammar and finite call-result domains, retain all
consistent symbolic alternatives, refine only by constraints valid for every
accepted program, and avoid revisiting the same finite symbolic state. A
fair traversal terminates and is complete for those hypotheses when the
checks are total. If result domains are infinite, merely adding fuel does
not create that theorem. Fuel exhaustion is unknown, and unbounded fair
refinement may be a semidecision strategy rather than a terminating one.

Contract equations such as $f(v)=w$ may constrain a recursive oracle on $v$
when every accepted candidate must satisfy that equation. Arbitrary outputs
invented for unobserved calls have a different status. Keep this distinction
in the constraint representation so that a convenient synthesis assumption
cannot silently become a pruning fact.

\section{R4: Indexed families, motives, and transport}
\label{sec:r4}

\question{R4.1}{Which motive enumeration can be complete and terminating?}

\answer A finite occurrence-abstraction and generalization grammar can be
exhaustively searched. Its completeness is relative to that grammar, and it
is not automatically closed under all accumulator generalizations.

Fix the recursor, a finite set of eligible maximal index occurrences, a
finite collection of context parameters, and a size bound on product,
function, and dependent-pair extensions. Enumerate subsets of occurrences
to abstract; compute dependency closure for each proposed reversion; form
the motive and ask Lean to type-check it. With $k$ eligible occurrences
there are at most $2^k$ raw subsets before typing and symmetry filtering.
Adding an unbounded number of auxiliary parameters destroys finiteness, so
those extensions must be bounded separately.

The first candidate should be Lean's ordinary motive, not an expensive
powerset search. Next try dependency-closed generalization of indices and
parameters actually implicated in a failed induction-hypothesis application.
Only then enter the bounded schema grammar of R2. Useful invariants can
require new types or relations absent from the target, so abstraction of
existing occurrences alone cannot provide unrestricted completeness.

For example, mapping a vector has the motive
\[
 M(n,xs)=\mathsf{Vec}(B,n).
\]
A function threading an accumulator indexed by the remaining length may need
$M(n,xs)=\Pi a:S(n),\,T(n,xs,a)$. This is not merely choosing a subset of
occurrences in the original result type. The common schema representation
should make the added generalized telescope explicit.

\question{R4.2}{How should propositional transports be represented and ranked?}

\answer Keep native Lean terms and exact types authoritative. Use certified
transport simplification as a secondary normalization layer, not an erased
setoid language that postpones all typing until the end.

Proof irrelevance identifies proofs of the same proposition. It does not
make propositionally equal indices definitionally equal. A term of
$\mathsf{Vec}(A,n+m)$ can require transport to inhabit
$\mathsf{Vec}(A,m+n)$ even when an arithmetic theorem proves the index
equality. The transport endpoints and their dependent family are therefore
part of the computational term's typing information.

Canonicalization can remove reflexive casts, compose adjacent casts, and
normalize index expressions using proved equalities. Store the proof needed
to justify each replacement. Heterogeneous equality is useful for relating
values at different types, but does not remove the obligation to produce a
term of the exact requested type. A setoid can be an auxiliary search
abstraction when its simulation and reconstruction obligations are explicit;
it is not a general substitute for Lean's dependent conversion checker.

Giving all casts zero cost can introduce unbounded families of superficially
free terms. Use a positive secondary syntax bound, or normalize within a
proved finite saturation discipline. Proof arguments may receive a small
ranking weight without being ignored for termination of enumeration. Hiding
casts in a displayed program is acceptable only when re-elaborating that
presentation reconstructs a checked term of the same intended type; otherwise
show the cast or an explicit local equality proof.

\question{R4.3}{How aggressively should definitions be unfolded during search?}

\answer Use staged, demand-driven unfolding. Neither full normalization of
all states nor a single shallow transparency regime is an adequate general
policy. Canonical demonstrates dependent inhabitation with definitional
equality; Aesop's discussion also illustrates the practical limitations of
head indexing through definitions. These are precedents, not evidence of a
universal optimal transparency schedule.\cite{canonical,aesop}

Begin with reducible weak-head normalization sufficient to expose the target
or provider head. On a rigid mismatch, selectively unfold definitions that
block a plausible application. Record the transparency level and unfolded
head in the cache key. A later, bounded fallback can use broader unfolding
rather than permanently treating a shallow mismatch as impossibility.

Normalize index arithmetic with checked lemmas where definitional reduction
is insufficient. Do not eagerly normalize a large open recursor merely to
compute an index key. Separate a failed attempt at one transparency setting
from a proof of non-unifiability. In particular, a search-level policy that
chooses not to unfold a definition is not the kernel proving that two
expressions cannot be convertible.

The right cost policy is empirical: count successful matches unlocked by
additional unfolding, normalization work, and repeated cache misses. Native
Lean unification should remain the final check rather than a home-grown
approximation of dependent definitional equality.

\question{R4.4}{Can Lean's motive computation be used as a library?}

\answer Yes, with version-pinned adapters and careful treatment of its
preconditions. The inspected Lean 4.34.0 source exposes
\code{Lean.Meta.mkRecursorAppPrefix}, whose documented role is to construct
a recursor prefix with the motive obtained by abstracting the major premise
and indices. It also exposes \code{Lean.MVarId.induction}, returning an array
of \code{InductionSubgoal} values with fields and a free-variable
substitution.\cite{leaninduction}

This gives a concrete integration path instead of implementing a second
motive compiler. The same source's \code{getMajorTypeIndices} requires
indices to be distinct free variables and checks dependency restrictions.
A major premise indexed by an expression such as $n+1$ therefore needs the
appropriate higher-level preparation; the low-level API is not a magic
solver for arbitrary index expressions.

Wrap each attempt in a complete search transaction. Report structured failure
categories such as nonvariable index, inadmissible elimination sort,
ill-typed generalized telescope, unresolved instance, or resource limit.
Memoize only under a key that includes the target, local context, recursor,
proposed generalization, and relevant assignment state. Resource exhaustion
must not become a permanent ``this motive is impossible'' entry.

Start with the high-level induction behavior already available in Lean,
then use the lower-level helper when controlled experiments show a benefit.
This preserves Lean's handling of dependent elimination and keeps maintenance
localized when the toolchain version changes.

\section{R5: General recursion from incomplete examples}
\label{sec:r5}

\question{R5.1}{Does a capability interpreter agree with the realized recursive term?}

\answer Agreement must be established by a correspondence theorem. For some
structural recursors it can be definitional; for well-founded recursion the
reliable general statement is propositional equality using the defining
equation and well-founded induction. A search-time interpreter is not
certified merely because its syntax resembles a recursive definition.

Represent a recursive problem by a state type $S$, a result family
$B:S\to\Type$, a well-founded relation $\prec$ on states, and a body with
access to a recursive capability
\[
 \mathsf{rec}_x:(y:S)\to(y\prec x)\to B(y).
\]
The state must include every recursive argument, including accumulators and
any tag needed for mutual recursion. The decreasing relation may depend on
only part of the state, but admissibility of every call must be proved for
the complete argument tuple.

\begin{proposition}[Interpreter--realizer agreement]\label{prop:recursion}
Suppose two total interpretations of a typed recursive body obey the same
recursive equations over a well-founded relation, and the body preserves
pointwise equality of its recursive-call results. Then the interpretations
agree at every state.
\end{proposition}
\begin{proof}
Use well-founded induction on $x$. Every recursive call made by the body at
$x$ has an argument $y\prec x$, where the induction hypothesis identifies
the two results. The body's equality-preservation property identifies the
result at $x$. This proves pointwise agreement everywhere.
\end{proof}

For a reflected body language, prove equality preservation by induction on
its constructors. If the proof instead uses equality of higher-order
recursive oracles, account for the function-extensionality principles used
under the selected Lean axiom profile. Do not silently assume that every
meta-level interpreter is extensional in the required sense.

A fuelled interpreter should return either a certified value or unknown.
Its theorem has the form
\[
 \mathsf{evalFuel}(p,x,n)=\mathsf{some}(v)
 \ \Longrightarrow\ \mathsf{realize}(p)(x)=v,
\]
not a claim that failure to return a value disproves the contract. A reduction
trace can reconstruct this equation from the realized function's equations.
That is often a more controllable certification path than asking kernel
normalization to unfold a large well-founded fixpoint directly.

During search, decrease proofs can themselves contain holes. An interpreter
must not treat a speculative recursive call as certified merely because its
measure appears smaller on one sample. Either establish the call's required
proof, carry it as a guarded obligation, or keep resulting deductions at the
heuristic level until the correspondence hypotheses hold.

\textbf{Implementation recommendation.} Introduce a small typed body language
only for recursive-call structure, keeping ordinary expressions and types
native to Lean. Realize surviving closed bodies through Lean's recursive
definition machinery, retain their equation lemmas, and verify the resulting
program and contract proof with the existing gate. The official reference
also supplies functional induction principles for supported recursive
definitions, which are useful for later contract proofs.\cite{leanmanual}
Compiled evaluation may rank or test candidates quickly, but a kernel-only
refutation policy must reconstruct evidence before deleting a branch on the
basis of an untrusted compiled result.

\question{R5.2}{Does angelic execution compose with top-down search?}

\answer Yes, if angelic choices are retained as correlated constraints rather
than prematurely committed values. The main loss relative to a bottom-up
version-space approach is that top-down refinement may repeatedly rediscover
similar residual behavior unless memoization and shared constraints are
provided. Burst is the relevant synthesis precedent, not a proof that any
top-down adaptation is complete.\cite{burst}

In a top-down body, an unknown recursive call produces a fresh symbolic
result tied to its typed argument tuple and current function candidate.
When two calls have equal arguments, their results must agree. If a contract
contains an equation determining that result, constrain it accordingly. If
not, derive a relation on possible results by backward propagation. Preserve
disjunctions when more than one result can explain the required output.

For example, seeing a residual equation $r_1+r_2=5$ does not justify assigning
$r_1=0,r_2=5$ and discarding the body if that assignment later fails. That is
only one explanation. Moreover, if the two symbols represent the same
recursive call, the residual must include $r_1=r_2$. Ignoring that correlation
admits spurious angelic explanations and can make refinement ineffective.

A practical hybrid uses top-down native typing to choose the recursion
scheme and body skeleton, with bottom-up enumeration of small closed
components in observed environments. The component store may group terms by
observed behavior as an index, while retaining alternatives when future
contexts can distinguish them. Eventually execute the complete candidate
without angelic choices and discharge the full contract. Unknown recursive
values help steer synthesis; they do not become new axioms.

\question{R5.3}{Is there a canonical small basis of recursion schemes with known coverage?}

\answer No verified measurement found in this investigation establishes such
a basis for Leant2's language, trust policy, and benchmarks. Scheme families
are useful engineering choices, but ``canonical'' and ``covers textbook
programs'' are not precise enough to imply a theorem or a percentage.

Use an ordered collection with separately reported coverage: ordinary
structural recursion including \code{Nat}; structural recursion with
parameters or tupled results; continuation-style accumulators; course-of-values
recursion; and well-founded recursion with a bounded grammar of measures and
lexicographic combinations. Add mutual and nested cases through explicit
Lean-compatible adapters when corresponding benchmark families justify the
cost. One general well-founded scheme may be semantically expressive, yet
far worse for search than several specialized schemes.

Fibonacci illustrates why scheme and carrier should not be evaluated
independently: course-of-values recursion realizes one presentation, while
ordinary natural-number recursion with a pair realizes another. List zip
can be expressed through recursion on one list with the other as a generalized
parameter. The best search representation is not determined by the surface
number of recursive arguments.

Evaluate the marginal gain of each scheme on family-held-out tasks, with
both known-carrier and invented-carrier tracks. Publish failures by cause:
no expression in the current grammar, body search interrupted, termination
proof unresolved, behavioral counterexample, or contract proof unresolved.
Without this separation, a scheme can appear to lack expressiveness when
the actual bottleneck is a missing library component or proof lemma.

\section{R6: Contracts beyond decidable observations}
\label{sec:r6}

\question{R6.1}{How should value holes and proof holes be interleaved?}

\answer Treat both as constraints, but distinguish cheap propagation from
expensive proof search. Waiting for every program hole to close is too late;
running arbitrary tactics on every partial state is too expensive.

The proposition that proof search is meaningless until its program holes
are filled is too strong. A proof goal such as $\mathord?n=3$ can determine
$\mathord?n$ by reflexivity and unification. A constructor proof or a
subtype invariant can similarly restrict an implementation choice. The
proper restriction is that a tactic attempt must preserve all shared
assignments transactionally and should run only when its expected benefit
justifies the work.

Use compositional verification conditions for recognized constructors,
recursors, and certified components. Cheap operations---unification,
constructor decomposition, reflexive equality, and small reflected
arithmetic---run eagerly. More expensive simplification, induction, and
library search are separate obligations with dependency watchers. Synquid's
refinement-based decomposition is a useful conceptual precedent, but Lean
propositions are richer than the decidable refinement fragment of that
system.\cite{synquid}

\subsubsection*{A complete local proof pattern: length-preserving map}

For $f:A\to B$, synthesize
\[
 F(xs):\{ys:\List(B)\mid \len(ys)=\len(xs)\}.
\]
The nil branch returns $(\nil,\mathsf{refl})$. In the cons branch, suppose
the recursive result is $(ys,h)$ with $h:\len(ys)=\len(xs)$. Return
\[
 \bigl(f(a)\cons ys,\ \mathsf{congrArg}(\mathsf{Nat.succ},h)\bigr),
\]
after the definitional reductions of list length. The proof obligations are
local and follow exactly the data construction. Projecting the first
component yields a map function with a length-preservation proof.

This pattern avoids asking a generic tactic to discover the entire induction
about an anonymous recursor after synthesis has finished. It also provides
constraints during construction: choosing a branch output of the wrong
length creates a local obligation that can often be refuted immediately.

A concrete prerequisite is visible in the repository: the current
\code{proofPortfolio} returns without trying tactics when the target contains
free variables. It cannot simply be reused unchanged for these local branch
obligations, which intentionally contain variables and an induction
hypothesis. Introduce a separate local-context proof service with explicit
scope and transactional semantics.\cite{leantcore}

A Lean-only analogue of finite liquid-style abstraction can start from a
bounded set of proposition templates involving lengths, membership, ordering,
and selected index equations. Infer and test candidate invariants, then
prove the selected ones with Lean. Without a decision procedure for the
chosen templates, this remains an incomplete proof-search discipline, not a
new complete refinement checker.

\question{R6.2}{Which properties are covered by induction, simplification, and arithmetic?}

\answer Define a sufficient, certificate-producing class rather than promise
a complete syntactic characterization of all theorems that a tactic might
prove.

Fix a structural recursion scheme, a proposed invariant $P$, an oriented
finite set of computation and rewrite lemmas, and a decidable residual
logic. A property belongs to the proposed automation class when every branch
verification condition, after applying the recursive hypotheses and the
specified rewrites, reduces to a formula in that residual logic and its
decision procedure proves it. This membership test is operational and
certificate-producing. It need not predict success without constructing the
verification conditions.

\begin{proposition}[Recursion-aligned verification]
Suppose each constructor branch of a structurally recursive program proves
its postcondition from the postconditions of all permitted recursive calls.
Then structural induction proves the postcondition for every input.
\end{proposition}
\begin{proof}
Apply the datatype's induction principle. Its hypotheses supply exactly the
postconditions assumed by the corresponding certified branch. Applying that
branch proof establishes the induction conclusion.
\end{proof}

Length preservation, affine size bounds, and many constructor-level
invariants fit a first implementation when their residual arithmetic is
Presburger arithmetic. Properties involving permutations, nested uses such
as $f(f(xs))=xs$, or nonlinear arithmetic may require auxiliary lemmas or a
different invariant. They are not automatically impossible; they are outside
the initial guaranteed tactic recipe until additional certified rules are
provided.

Keep a provenance record of the recursion scheme and the equations used.
It makes branch-condition generation cheap and avoids reverse-engineering
an arbitrary kernel term. Successful helper lemmas should be abstracted over
their contexts and stored as checked reusable components, subject to the
same trust profile as ordinary providers.

\question{R6.3}{What automated induction is usable in Lean now?}

\answer Lean 4.34 provides ordinary induction and the more targeted
\code{fun_induction} and \code{fun_cases} facilities. The reference describes
functional induction for supported non-mutually-recursive structural or
well-founded definitions, including generalization controls. These are
concrete building blocks already missing from the current small portfolio.
They do not constitute autonomous discovery of every needed induction or
lemma.\cite{leanmanual}

Combine a bounded choice of induction variable or synthesized function with
simplification by its equations, the local induction hypotheses, and
\code{omega} or other checked closers. \code{grind} is useful general-purpose
proof automation, but it should not be described as an automatic replacement
for choosing all induction principles and inventing strengthening lemmas.
Optional Aesop-style rule search and Canonical-style inhabitation supply
additional approaches.\cite{aesop,canonical}

A raw recursor expression does not necessarily have the named functional
induction theorem expected by a high-level tactic. Preserve the original
recursive equations and realize promising programs as definitions when
using functional induction, or apply the relevant structural induction
principle directly. This is another reason to retain synthesis provenance.

The current Lean Hammer work combines premise selection with external
reasoning and checked proof reconstruction. It is relevant to a later
optional proof lane, not a core-only, dependency-free induction solution.
Its measurements concern theorem proving, not automatic certification of
Leant2's invented recursive programs.\cite{leanhammer}

Use Lean resource limits and interrupt-aware wrappers for every tactic call.
A library interface with a fuel or heartbeat limit provides controlled
work, not necessarily a perfectly predictable wall-clock duration. Failed
or interrupted attempts leave the candidate unproved unless a separate
counterexample or contradiction has been certified.

\question{R6.4}{Is there a principled policy for proof debt?}

\answer There are useful decision-theoretic heuristics, but the coupled,
state-changing proof obligations here do not automatically satisfy the
assumptions of a classical independent-arm bandit theorem.

For comparison, suppose two independent candidate tests have fixed costs
$c_1,c_2$ and rejection probabilities $p_1,p_2$, and testing stops at the
first rejection. Testing 1 then 2 costs in expectation
$c_1+(1-p_1)c_2$; the reverse order costs $c_2+(1-p_2)c_1$. The first order is
better exactly when $p_1/c_1\geq p_2/c_2$. This justifies a simple
rejection-per-cost heuristic for suitable tests. Proof attempts differ:
success may discharge several coupled obligations, failure is often
inconclusive, and both outcomes can change subsequent search.

Use a deterministic tiered policy. Eager propagation pays very small local
obligations. A bounded medium tier attempts simplification and arithmetic
when their data dependencies stabilize. An expensive tier considers induction
and retrieval after the candidate has survived known counterexamples and
its recursive structure is stable. Reserve a fixed share of work for older
pending obligations so that promising but difficult states are not starved.

Limit \emph{active} proof debt by suspending states, not by declaring them
false. Charge proof work to the same logical ledger as construction and
evaluation. Learn cost estimates for ordering if useful, while keeping the
fairness rule and immutable ranking objective unchanged. Before acceptance,
every required proof debt must be discharged. The exact thresholds are
benchmark parameters to calibrate for a release, not quantities already
known to be optimal.

\section{R7: Component retrieval at library scale}
\label{sec:r7}

\question{R7.1}{What type abstraction should support dependent reachability?}

\answer Use a sound typed AND-hypergraph abstraction with refinement, not only
a graph of result-type heads. Preserve a native-Lean checking boundary.
Hoogle+'s type-guided abstraction refinement is the closest directly useful
precedent, but Lean adds genuinely dependent indices, universe constraints,
and a richer instance mechanism.\cite{tygar}

A component with several required arguments is an AND-transition: all its
preconditions must be realizable with compatible substitutions. A simple
path from some available type head to the goal head can ignore an impossible
second argument. It can also forget that the same type variable occurs in
several positions. Therefore abstract nodes should retain normalized head
structure, a bounded pattern of shared type variables, selected index
predicates, universe relationships, and explicit dictionary requirements.
The abstraction may forget additional constraints as long as it
\emph{over-approximates} concrete applicability.

Lean's ordinary local variables are reusable. If a Petri-net representation
is adopted, it must represent contraction and weakening appropriately;
consuming an input token linearly without those rules would discard valid
programs. Likewise, the abstract representation of a dependent component
must include every concrete instantiation promised by the retrieval policy.
Otherwise abstract unreachability is not a sound reason to omit it.

When a concrete path fails Lean type checking, inspect the mismatch and add
a useful predicate to the abstraction: an index equality, a relationship
between repeated type parameters, or a dictionary requirement. Merely
checking each proposed term by unification is not itself abstraction
\emph{refinement}; the next abstract search must actually use the newly
learned distinction.

A path-length cutoff of three can be an effective fast lane. It is an
incomplete heuristic unless the declared search grammar itself has that
composition bound. Maintain a fair broader fallback, or state the cutoff in
any exhaustion result. Use indexing to shortlist plausible components;
final native typing decides whether an actual application is legal.

\question{R7.2}{Has type-guided abstraction refinement handled type classes?}

\answer Yes. Hoogle+'s implementation translates type classes into explicit
dictionary types: class operations consume dictionaries, and instances
construct them. Its discussion also identifies limitations, including
unsupported higher-kinded type variables such as the parameter in
\code{Monad m}. Thus the starting point is an existing construction to adapt,
not an entirely unexplored type-class problem.\cite{tygar}

The corresponding Lean representation should use actual instance arguments.
An instance constraint is a typed obligation, not merely a Boolean label
saying that a class name occurs. Resolving it may assign metavariables,
depend on local instances, or fail because parameters are still unknown.
Run speculative instance search inside the branch transaction and suspend
underdetermined requests rather than caching them as permanently impossible.

A cache key should include the full normalized class application, universe
instantiation, relevant local context and local instances, environment
fingerprint, and assignment dependencies. The proposal's shorthand
``(class, type)'' key is insufficient for classes with several parameters or
dependent arguments. Instance dictionaries can also affect behavior: two
valid structures on the same carrier type may interpret operations
differently. Behavioral observations and compiled-evaluation caches must
retain which dictionary was chosen.

Implement dictionary-aware filtering before ambitious dependent refinement.
Measure whether it reduces the number of failed component applications
without causing missed successes under the declared abstraction. Gradually
retain more index and universe information only where traces show that
coarse abstraction generates excessive spurious paths.

\question{R7.3}{Can Lean's suggestion interface accept data goals and examples?}

\answer Its type accepts an arbitrary goal metavariable, and its configuration
already has an optional textual hint. That is a useful integration point,
but it is not a guarantee that an installed selector understands data
synthesis or examples.

In the inspected Lean 4.34.0 source,
\begin{lstlisting}
abbrev Selector := MVarId -> Config -> MetaM (Array Suggestion)
-- Selected Config fields:
maxSuggestions : Nat
caller         : Option String
filter         : Name -> MetaM Bool
hint           : Option String
\end{lstlisting}
The source explicitly says that Lean does not provide a default library
suggestion engine. The hint is arbitrary and a selector may ignore it.
There is no required structured input--output-example field in this
interface.\cite{leansuggestions}

A Leant2 adapter can pass a data goal, identify itself through the caller
field, and serialize a compact observation summary into the hint for a
selector designed to consume it. For a reliable local default, use the
native provider index and abstract component graph directly. A specialized
selector can add typed behavioral features without changing the Lean
interface, but exact application checking remains necessary downstream.

A training corpus can be built from suitable Std or Mathlib definitions:
retain their type and approved context, hide the body, extract used
components as labels, and generate well-typed observations when the function
is executable. This is a dataset proposal, not an existing measured
Leant2 corpus. Hide aliases, generated equations, and directly exposing
helper declarations where they would leak the answer. Split by module,
chronology, and near-duplicate algorithm family, not merely by randomly
selected definitions. Keep instance choices and the environment snapshot
fixed. General lemmas may remain available when the benchmark is explicitly
about realistic library use, but disclose that availability.

The Lean Hammer paper describes retrieval over theorem-proving states and
library signatures. Its strong premise-selection results do not establish
relevance accuracy for data goals constrained by examples.\cite{leanhammer}
Use a separate held-out evaluation for the latter. Network-backed selectors
must remain an explicit optional extension; a core-only default should not
silently send user goals or environment contents to a service.

\section{R8: Ranking and candidate equivalence}
\label{sec:r8}

\question{R8.1}{What is a defensible ranking objective?}

\answer Minimum description cost under a fixed grammar is defensible as a
stated prior, not as a theorem that the first answer is the user's intended
program. A useful objective can combine code, proof, execution, and
presentation costs:
\begin{equation}\label{eq:ranking}
 J(p,\pi)=-\log P_{\theta}(p)
 +\lambda_{\pi}\,\mathsf{proofCost}(\pi)
 +\lambda_t\,\mathsf{runCost}(p)
 +\lambda_h\,\mathsf{presentationCost}(p).
\end{equation}
For optimality claims, fix $\theta$, the weights, the cost definitions, and
the candidate grammar for the run. An estimated runtime term is a ranking
feature; it is not a proved complexity bound unless separately certified.
Probabilistic grammars and learned enumeration costs are established ideas,
but their adaptation to user-preferred Lean inhabitants requires
measurement.\cite{probe}

A grammar distribution is often over derivations, not programs. Multiple
derivations of the same term can distort both probability and size.
Normalize construction paths or define program cost as the minimum cost of
its permitted derivations. The implementation must match the chosen
mathematical objective before claiming that a cost-ordered search is optimal.

Ignoring an explicit argument is not always suspicious: constants and
projections are frequently correct. Input-use penalties should therefore be
soft, contract-aware features rather than a semantic rejection rule.
Similarly, syntactic occurrence is only an over-approximation of influence;
$x-x$ mentions $x$ without depending on it extensionally. Exact information
flow is not obtained merely by counting variable occurrences.

\begin{proposition}[A valid stopping certificate]\label{prop:stopping}
For a fixed objective $J$, suppose every unexplored or uncertified candidate
is covered by a frontier state with a valid lower bound $L(s)$ on its final
cost. If a certified candidate has cost $J^*$ and
$J^*\leq\inf_s L(s)$, then no covered candidate has smaller cost.
\end{proposition}
\begin{proof}
Every covered candidate has cost at least its state's lower bound, which is
at least $J^*$. Exhausted parts of the search contain no unexamined candidate.
\end{proof}

The frontier must include deferred proof obligations, suspended carrier
choices, and candidates discovered by any auxiliary lane. Adaptive search
priorities are not these lower bounds. A margin-based stopping rule gives
only the corresponding approximation guarantee; it does not establish
semantic uniqueness or intendedness.

Show a small diverse set of certified candidates, their important behavioral
differences, and a distinguishing input when one is found. For a ranking
benchmark, compare reference rank and family-level alternatives, but do not
treat one reference implementation as the unique acceptable program. Human
preference claims require an actual user study rather than inference from
shorter syntax.

\question{R8.2}{Which equivalence is cheap and meaningful for dependent terms?}

\answer Use a hierarchy and name the level achieved. No single cheap test
serves kernel equality, extensional equality, and useful behavioral diversity.

First canonicalize bound-variable names and approved syntactic normal forms.
Then use Lean definitional equality when economical. A separately checked
propositional equality can establish more identifications; function equality
or heterogeneous equality requires its corresponding proof and trust
assumptions. Observational agreement on generated inputs is a useful final
\emph{bucket}, but is not generally an equality proof.

For dependent terms, retain exact types and any transport witnesses. Proof
irrelevance can collapse proofs of the same proposition; it does not collapse
arbitrary values at propositionally related types. Type formers such as
\code{fun _ => Empty} and \code{fun _ => Unit} may have similar syntactic
shapes while being meaningfully different. Grouping them for presentation
must not be called semantic deduplication.

Finite observational quotienting is safe for destructive search pruning only
when it is a congruence for every future context admitted by the search and
all relevant environments. Agreement on the current examples rarely gives
that guarantee for open, higher-order, or dependent terms. Otherwise retain
multiple representatives lazily and split a bucket when a new observation
separates them. The elementary counterexample is the identity function and
the constant-zero function: they agree on input zero and differ on input
one.

Contract inputs alone often cannot separate any accepted candidates, because
all candidates already satisfy them. Add fresh well-typed probes and, in a
proved parametric fragment, the canonical position inputs from R2.2.
Report ``same on these probes'' unless a stronger equality certificate is
available. Printed strings should be presentation artifacts, not the primary
identity of terms or constraints.

\section{R9: Learned guidance under a kernel gate}
\label{sec:r9}

\question{R9.1}{How accurate are models at proposing carriers or helpers?}

\answer The inspected primary literature does not provide a verified number
for the exact question: Lean types plus sparse examples, invention of a
useful carrier or helper, successful realization and proof, and a short
end-to-end interactive budget. Theorem-proving and premise-retrieval results
cannot fill this gap by extrapolation.\cite{leanhammer,canonical}
This is a measurement question for which a focused experiment is more useful
than an unsupported percentage.

Separate proposal quality into stages. Does the suggested type elaborate?
Does it admit a body in the same bounded grammar used by the baseline?
Does the realized program satisfy the observations? Is the requested full
contract proved? Does the entire run, including model inference, meet the
budget? Report all stages, not just syntactically plausible suggestions.

Use paired comparisons against the same grammar with no model, a fixed
handwritten carrier list, and collision-directed summary invention. Evaluate
both known-carrier and invented-carrier tracks. Keep the model's permitted
components, examples, and axiom profile identical to those of the baseline.
Report top-$k$ proposal success and end-to-end verified success separately,
together with elaboration cost, proof cost, inference latency, hardware, and
resource use. A model that proposes the right carrier but consumes the
entire interaction budget has not achieved the requested end-to-end result.

Hold out algorithm families and close variants to reduce leakage from
textbook solutions. A carrier is not wrong merely because it is presented
by an isomorphic type, but acceptable equivalence and cost criteria must be
defined in advance. Count a proposed helper as useful only after its
implementation and required proof are obtained, or report it explicitly as
an unresolved sketch. Preserve failed proposals as well as successful ones
in the experimental record.

The best first insertion point is one bounded batch of schema proposals at
query preparation or a deterministic decision epoch, not a network call
inside every node expansion. The report recommends this as an experiment,
not as evidence that an optional model is already a net performance win.

\question{R9.2}{How can learned contributions be deterministic and auditable?}

\answer Make search replay consume immutable, typed proposal records.
Temperature zero and a fixed random seed alone are not a sufficient
specification of byte-for-byte reproducible inference across implementations
and hardware. Recording exact proposals separates proof replay from model
regeneration.

A receipt should identify the environment and toolchain, target and contract,
local telescope, grammar version, axiom profile, fixed ranking objective,
logical schedule, proposal bytes and order, checked terms, and certificates.
If model regeneration is itself a research objective, also identify model,
tokenizer, inference implementation, and relevant runtime configuration.
These are provenance requirements, not additional logical axioms. A signature
or digest protects provenance integrity; the kernel check establishes the
meaning of a program and its proof.

There is an important security boundary beyond logical soundness. A model
must not supply arbitrary Lean commands, elaborator extensions, or tactics
for execution inside the trusted process. Such code can have effects before
the final kernel gate runs. Accept a restricted typed abstract syntax or
strictly parsed term proposal over approved constants, and isolate any
untrusted elaboration that remains. The kernel gate does not undo I/O or
other effects produced by executing a proposal generator.

If learned suggestions only reorder a fixed grammar, the enumeration
argument remains applicable with a fair baseline lane. If they introduce
new carriers or helpers outside it, create an explicit grammar-expansion
epoch and update the meaning of exhaustion and ranking certificates. Never
claim both ``the model changes only order'' and ``the model can introduce
arbitrary new types'' without reconciling those two interfaces.

Finally, replaying an accepted proof establishes acceptance, not the absence
of a better candidate or the correctness of every prior pruning action.
Those stronger claims require the relevant search frontier and pruning
certificates. Keep the default implementation entirely model-free; optional
proposals can improve discovery without becoming part of the logical trust
base or making the ordinary core depend on service availability.

\section{A concrete Lean-native implementation plan}
\label{sec:implementation}

\subsection{Changes grounded in the inspected source}

The following table distinguishes integration points that exist from proposed
new services. Paths without a ``proposed'' qualification are present in the
inspected snapshot. This is a source-based plan, not a patch claimed to have
been built or benchmarked.\cite{leantrepo,leantcore,leantgate}

\begin{longtable}{@{}L{52mm}L{103mm}@{}}
\toprule
Location or boundary & Recommended change and its reason\\
\midrule
\endhead
\path{Leant2/Search/Core.lean}: alternatives and obligations & Extract a first-class state and queue policy while preserving the current depth-first order for regression comparison. Inventory every mutable reference for rollback behavior.\\
\path{Core.lean}: residual predicates and blockers & Give each observation its own status, dependencies, and branch-valid residual cache. Keep kernel-backed checking as the initial evaluator.\\
\path{Core.lean}: proof portfolio & Add a separate transactional proof service for local contexts. Its goals may contain locals and induction hypotheses, unlike the current closed-goal portfolio.\\
\path{Leant2/Accept/Gate.lean} & Retain synchronous kernel checking and axiom auditing. Make the binding between the requested contract and the supplied proof type explicit in the acceptance API.\\
Proposed constraint service & Track transitive expression and universe dependencies, delayed assignments, typed observation closures, and wakeups.\\
Proposed schema service & Enumerate bounded carriers, generalized motives, summary features, and recursive-call schemes through one typed interface.\\
Lean 4.34 induction adapter & Reuse \code{MVarId.induction} and \code{mkRecursorAppPrefix}; preserve returned substitutions and provenance.\\
Lean suggestion adapter & Consume the \code{Selector} API or the local provider graph; include examples only through a selector that explicitly understands them.\\
Proposed receipt/replay service & Record immutable query, grammar, schedule, proposal, certificate, and trust-profile data. Distinguish an interrupted prefix from exhausted search.\\
\bottomrule
\end{longtable}

The gate currently checks a supplied proof against its supplied proof type.
With a larger proposal ecosystem, require the boundary to verify that this
type is the requested $C(p)$, or a specified equivalent query result type,
for the exact accepted program. This is interface hardening suggested by
the function signature, not a reproduced erroneous acceptance in the current
frontend. A proof of an unrelated true proposition must not be relabeled as
a proof of the user's contract.\cite{leantgate}

\subsection{A state-expansion protocol}

The following is algorithmic pseudocode, not a compilable Lean API sketch.
It describes ordering and certificate boundaries rather than new trusted
reduction rules.

\begin{minipage}{\linewidth}
\begin{lstlisting}
expand(state):
    restore_complete_branch_state(state)
    saturate_bounded_cheap_typing_constraints()
    wake_observations_whose_dependencies_changed()

    for observation in ready_observations:
        result := evaluate_with_evidence_or_unknown(observation)
        if result contains a stable_refutation_certificate:
            record_refutation_and_stop_this_branch()

    if all_program_and_proof_holes_are_closed:
        check_exact_query_binding()
        return acceptance_gate(program, contract_proof, profile)

    if a local proof obligation is due under the fixed work policy:
        run_bounded_transactional_proof_service()
    else:
        decision := choose_ready_constraint_or_schema_hole()
        proposals := fair_merge(grammar_choices, recorded_guidance)
        for proposal in proposals:
            child := apply_in_fresh_transaction(proposal)
            if child is well_typed_or_has_valid_pending_constraints:
                enqueue(child, immutable_cost, adaptive_priority)
\end{lstlisting}
\end{minipage}

Every unknown branch remains eligible under the declared fair schedule unless
it is explicitly outside the current finite grammar. Observation and proof
services return typed outcomes, not a Boolean that conflates failure,
unsupported syntax, resource exhaustion, and refutation. This discipline is
more important than whether the first queue implementation is a heap, stack,
or a deterministic portfolio of both.

\subsection{Schema invention as a bounded subproblem}

A schema proposal should contain its carrier or generalized telescope, seed,
step holes, finishing hole, recursive-call permissions, and invariant holes.
For a fold, the input and output interface is
\[
 (A,T,C,b)\longmapsto
 (K,z,s,d,\mathcal V),
\]
where $\mathcal V$ contains the typing, behavioral, and invariant obligations
needed to certify the schema. A carrier suggestion without a realizable step
and finish is an intermediate hypothesis, not a successful synthesis.

Store a version-space of bounded feature proposals rather than committing
to the first collision separator. If one feature can be computed on whole
inputs but cannot be updated compositionally in the current step grammar,
try the next feature set. Helper invention must obey the same rule: a helper
with an underdetermined observation table remains a synthesis problem until
its implementation and required properties have been discharged.

This common interface permits a small default library of useful schemas,
collision-directed suggestions, and optional recorded model proposals to
coexist. It also makes ablations meaningful: all proposal sources feed the
same realization, evaluation, and certification pipeline.

\subsection{Dependency-ordered implementation sequence}

\paragraph{Phase A: preserve behavior while exposing guarantees.}
Refactor state management without changing the old enumeration order. Add
structured outcomes, a logical work ledger, exact query binding at acceptance,
and deterministic receipts. Test sibling-branch rollback and replay before
changing heuristics. This phase establishes the conditions under which later
measurements mean what they claim.

\paragraph{Phase B: obtain more information from the existing evaluator.}
Split recognized observations through checked rules, track dependencies per
observation, and add counterexample reuse. Do not first implement an entire
replacement kernel. Measure avoided reevaluations and the number of
observations that remain opaque or unsupported.

\paragraph{Phase C: add high-value native reasoning.}
Permit local-context proof obligations, integrate source-verified induction
adapters, and add demand-driven arithmetic and provider lookup. Introduce a
small number of typed branch and recursion schemas with explicit bounds.
This addresses gaps that do not require a breakthrough in carrier invention.

\paragraph{Phase D: bounded summary and motive invention.}
Add backtrackable feature grammars, sample incompatibility certificates, and
schema-specific invariants. Implement the carrier-given track first so that
failures of body synthesis are not misattributed to invention. Add true
abstraction refinement to data-component retrieval when spurious paths become
a measured bottleneck.

\paragraph{Phase E: certified fragment evaluation and richer recursion.}
Only after residual traces identify the expensive fragment, implement its
proof-producing evaluator. Add capability-based course-of-values or
well-founded recursion with an explicit interpreter--realizer theorem and
retained equations. Unsupported cases continue through slower native paths
or remain unknown, never false by default.

\paragraph{Phase F: learning and preference experiments.}
With reproducible traces and stable baselines in place, compare learned
proposals and ranking policies. Keep model-free behavior available and avoid
turning one optional optimization into a prerequisite for ordinary synthesis.

These phases have logical dependencies, not promised delivery durations.
The order differs from a plan that treats carrier invention as the prerequisite
for every other gain. Correct local proof handling, per-observation evaluation,
and native indexed induction can improve the system before a general carrier
inventor is available.

\section{Evaluation, reproducibility, and explicit limits}
\label{sec:evaluation}

\subsection{A benchmark that distinguishes the bottlenecks}

Retain the existing regression suites, but add tasks partitioned by the
capability needed: conditionals and arithmetic; ordinary and accumulating
recursion; multiple recursive results; indexed families and transport;
universal contracts; large-library composition; and underdetermined
behavioral specifications. Record the allowed providers, trust profile,
recursion schemes, type grammar, and example sets for every task.

Use paired tracks for carrier given versus invented, reference implementation
available versus examples only, and trace-complete versus incomplete examples.
For dependent tasks, include hand-checked target motives as diagnostic data,
without giving them to the invention track. A supplied motive that still
fails isolates body search, transport, or proof automation as the obstacle.

The main performance outcome is the number of \emph{verified} solutions at
specified wall-clock budgets, accompanied by logical work counts. Report
unrefuted/unproved candidates separately. Also report time to first verified
candidate, verification cost, peak retained frontier, residual evaluations,
cache validity checks, and the rank of an appropriate reference solution.
Avoid counting a finite-test survivor as a universally correct program.

Evaluation should include complete-cost accounting. A retrieval lane must
include index construction or explicitly report amortized warm-cache results.
A model lane includes inference and elaboration. A compiled evaluator lane
includes whatever evidence reconstruction is required by its trust policy.
Do not compare a kernel-only baseline with an unverified execution path as
if their guarantees were identical.

\subsection{Ablations and adversarial regression tests}

The most informative ablations disable one mechanism at a time: per-observation
invalidation, dependency-aware ordering, carrier lifting, local proof
obligations, indexed induction, abstract retrieval refinement, and learned
proposals. Use identical query snapshots and deterministic logical budgets
for algorithmic comparisons; use repeated paired runs for wall-clock effects.

Add adversarial cases specifically for the new invariants. Backtrack after a
metavariable assignment and revisit an identically printed residual. Introduce
local instances with different behavior. Couple two typing-independent data
holes by one observation. Use an intentionally non-lawful Boolean equality
instance. Interrupt a proof attempt that would eventually succeed. Offer two
terms equal on training inputs but different on a fresh one. Propose an
ill-scoped dependent closure. Introduce a speculative recursive call whose
decrease proof is still open. Every case should have an expected outcome
class, not just an expected printed candidate.

For ranking certificates, put the eventual cheaper program in a deferred
carrier or proof lane. The engine must not stop with an optimality claim
while that lane has a smaller valid lower bound. For learned guidance,
include a proposal that lies outside the declared grammar and verify that
its admission changes the grammar epoch rather than silently invalidating
an exhaustion certificate.

\subsection{The included finite experiment}

The archive contains a dependency-free Python script and its JSON output.
Its role is to make several small counterexamples and finite counts
reproducible. It is not a prototype of the Leant2 search engine.

\begin{center}
\begin{tabular}{@{}L{91mm}r@{}}
\toprule
Check & Result\\
\midrule
Distinct observations for the four-state lower bound & 9\\
One-state machines exhaustively rejected & 2\\
Two-state machines exhaustively rejected & 128\\
Three-state machines exhaustively rejected & 17,496\\
Binary inputs checked against the four-state realization & 2,047\\
Ternary inputs checked for the direct reverse fold & 3,280\\
\bottomrule
\end{tabular}
\end{center}

The script additionally checks elementary witnesses for nontransitive partial-row
compatibility, the invalid shared-goal cost sum, finite-probe nonequivalence,
branch-blind cache invalidity, and transitive dependency coupling. These are
small mathematical regression examples, not evidence that the corresponding
bugs have occurred in Leant2. The reverse equality and the minimal four-state
claim have proofs in Section~\ref{sec:r2}; finite testing is not substituted
for those proofs.

Run the experiment from the archive directory with
\begin{lstlisting}
python3 experiments/check_claims.py --output experiments/results.json
\end{lstlisting}

No Lean compiler was available in the execution environment for this report.
Consequently no Lean snippet was compiled, no proposed tactic adapter was
executed, and no repository benchmark was independently reproduced. The API
claims come from reading the pinned source and versioned documentation.
The algorithmic pseudocode is deliberately not presented as tested Lean code.

\subsection{What remains genuinely uncertain}

The largest empirical uncertainty is not whether more carrier expressions
can be enumerated, but whether collision-driven lifting sufficiently reduces
the joint search for features, steps, finishes, and proofs on realistic
queries. The second is whether typed residual reuse remains cheaper than
reduction after accounting for dependency bookkeeping and branching. The
third is whether generalization and local verification conditions cover a
large useful fraction of dependent programs without uncontrolled branching.

Model accuracy and human ranking preferences remain unmeasured for the exact
Leant2 workload. A strong result in theorem proving or examples-only synthesis
is evidence of a technique's potential, not a substitute for those
measurements. None of these uncertainties prevents implementing the bounded
algorithms and testing the hypotheses. They instead determine which claims
should be made as design proposals and which require new experiments.

\section{Conclusion}

The expert questions do not all have the same status. Some have concrete
answers in existing work or current Lean APIs: higher-order live examples,
dictionary-aware type-guided abstraction, assignment-triggered constraints,
and reusable motive construction. Some have precise restricted solutions:
minimal sample-consistent finite machines, bounded motive enumeration,
recursion-aligned proof obligations, and abstract-distance search lower
bounds. Others are empirical design questions, especially efficient carrier
invention, proof scheduling, learned suggestions, and user-preferred ranking.

The strongest unifying recommendation is to preserve native Lean typing and
make uncertainty and evidence explicit at every boundary. A sound gate,
a conservative partial evaluator, a fair bounded enumerator, an immutable
ranking objective, and a deterministic replay mechanism solve different
problems. Conflating them creates attractive but invalid shortcuts.

A practical next version of Leant2 can make substantial progress without
assuming a general solution to dependent program synthesis. Start with
explicit states, observation-level dependencies, local proof obligations,
and Lean's own induction machinery. Then add bounded, backtrackable summary
invention and certified recursive realization. This creates a measurable
path from the current synthesizer to stronger programs and contracts, while
retaining a small and intelligible logical trust boundary.

\clearpage
\appendix
\section{Question-by-question coverage map}

Every enumerated question in the nine source research areas has an explicit
answer. R8 and R9 share the source file \path{12-ranking.tex}; they are
separate areas in this report. The source's consolidated expert table is a
priority summary, not an additional set of questions.\cite{leantproposal}

\begin{longtable}{@{}L{15mm}L{121mm}r@{}}
\toprule
ID & Subject of the answer & Page\\
\midrule
\endhead
R1.1 & Shared metavariables, independence, and cross-branch memoization & \pageref{q:R1.1}\\
R1.2 & Admissible lower bounds with coupled goals & \pageref{q:R1.2}\\
R1.3 & Just-in-time learning and ranking guarantees & \pageref{q:R1.3}\\
R1.4 & Wall-clock portability and deterministic replay & \pageref{q:R1.4}\\
R2.1 & Minimal carriers from finite examples & \pageref{q:R2.1}\\
R2.2 & Polymorphic alphabets and canonical position inputs & \pageref{q:R2.2}\\
R2.3 & Backward fusion and auxiliary-summary invention & \pageref{q:R2.3}\\
R2.4 & Motives for primitive recursion with parameters & \pageref{q:R2.4}\\
R2.5 & Generalization and lemma-invention heuristics & \pageref{q:R2.5}\\
R3.1 & Higher-order and dependent live evaluation & \pageref{q:R3.1}\\
R3.2 & Incremental normalization and rollback invalidation & \pageref{q:R3.2}\\
R3.3 & A conservative evaluator and the role of Lean4Lean & \pageref{q:R3.3}\\
R3.4 & Angelic top-down termination and completeness & \pageref{q:R3.4}\\
R4.1 & Finite motive enumeration and generalization closure & \pageref{q:R4.1}\\
R4.2 & Canonical transport handling & \pageref{q:R4.2}\\
R4.3 & Definitional unfolding policies & \pageref{q:R4.3}\\
R4.4 & Lean motive-computation library interfaces & \pageref{q:R4.4}\\
R5.1 & Interpreter--recursor correspondence and evaluation & \pageref{q:R5.1}\\
R5.2 & Angelic execution in top-down synthesis & \pageref{q:R5.2}\\
R5.3 & A small recursion basis and measured coverage & \pageref{q:R5.3}\\
R6.1 & Interleaving implementation and proof holes & \pageref{q:R6.1}\\
R6.2 & A certifiable class of recursion-aligned properties & \pageref{q:R6.2}\\
R6.3 & Current Lean induction and proof automation & \pageref{q:R6.3}\\
R6.4 & Scheduling proof debt & \pageref{q:R6.4}\\
R7.1 & Dependent type abstraction for component reachability & \pageref{q:R7.1}\\
R7.2 & Type classes in type-guided abstraction refinement & \pageref{q:R7.2}\\
R7.3 & Data-oriented premise selection and training data & \pageref{q:R7.3}\\
R8.1 & Formal ranking objectives and stopping certificates & \pageref{q:R8.1}\\
R8.2 & Meaningful candidate equivalence & \pageref{q:R8.2}\\
R9.1 & Accuracy of learned carrier and helper proposals & \pageref{q:R9.1}\\
R9.2 & Deterministic and auditable learned guidance & \pageref{q:R9.2}\\
\bottomrule
\end{longtable}

\clearpage
\section{Source audit and interpretation of prior work}

\begin{longtable}{@{}L{40mm}L{115mm}@{}}
\toprule
Source & What was used, and what was not inferred\\
\midrule
\endhead
Leant2 proposal and code & All nine question areas and the expert summary; inspected construction, residual, proof-portfolio, and acceptance code. No independent performance reproduction.\\
Canonical-min (2026) & Assignment-triggered suspension and a small dependent inhabitation/unification implementation. Not a claim of drop-in full Lean compatibility.\\
Smyth (2020) & Higher-order examples in the implementation, live bidirectional evaluation, and removal of trace-completeness requirements. Not a certified evaluator for Lean's dependent kernel.\\
AutoLifter (2024 revision) & Auxiliary-summary and combinator synthesis; its own restrictions on minimality and greedy search. Not minimal Lean carrier types inferred uniquely from arbitrary finite samples.\\
Hoogle+ (2020) & Type-guided abstraction refinement and explicit dictionary encoding. Not unrestricted support for Lean dependency, universes, or all higher-kinded classes.\\
Lean 4.34 source/manual & Exact induction helpers, functional induction facilities, and the suggestion interface including its hint field and lack of a built-in selector engine.\\
Lean4Lean (2025 revision) & An external Lean typechecker and verification of parts of the checker/metatheory. Not a complete proof for every rule of a newly proposed partial evaluator.\\
Lean Hammer (2026 revision) & Current premise-selection and proof-reconstruction architecture. Not an empirical estimate of carrier invention or data-synthesis accuracy.\\
\bottomrule
\end{longtable}

\addcontentsline{toc}{section}{References}
\begin{thebibliography}{99}
\small\raggedright
\setlength{\itemsep}{1pt}
\setlength{\parsep}{0pt}
\setlength{\parskip}{0pt}

\bibitem{leantrepo}
Vladimir Reshetnikov. \emph{Leant2 repository and README}. Inspected snapshot
\code{784664ff34e819422510a4d470ab07fe48bb736b}, September 18, 2026.
\url{https://github.com/VladimirReshetnikov/Leant2/tree/784664ff34e819422510a4d470ab07fe48bb736b}.

\bibitem{leantproposal}
Leant2 project. \emph{Leant2: Design for Further Improvements}. September 18,
2026. Sections R1--R9 and the consolidated expert-question table; same pinned
snapshot as [1]. The proposal's original measurements refer to an earlier
commit and are not reproduced in this report.
\url{https://github.com/VladimirReshetnikov/Leant2/tree/784664ff34e819422510a4d470ab07fe48bb736b/docs/proposals/11-further-improvements}.

\bibitem{leantcore}
Leant2 project. \emph{Construction search},
\path{Leant2/Search/Core.lean}. Pinned snapshot as [1]; especially
\code{alternative}, \code{instantiatePartial}, \code{mkResidual},
\code{kernelDecide}, \code{evalResidual}, \code{partialRefute}, and
\code{proofPortfolio}.
\url{https://github.com/VladimirReshetnikov/Leant2/blob/784664ff34e819422510a4d470ab07fe48bb736b/Leant2/Search/Core.lean}.

\bibitem{leantgate}
Leant2 project. \emph{The acceptance gate},
\path{Leant2/Accept/Gate.lean}. Pinned snapshot as [1].
\url{https://github.com/VladimirReshetnikov/Leant2/blob/784664ff34e819422510a4d470ab07fe48bb736b/Leant2/Accept/Gate.lean}.

\bibitem{aesop}
Jannis Limperg and Asta Halkj\ae r From.
Aesop: White-Box Best-First Proof Search for Lean.
\emph{CPP}, 2023. DOI: \href{https://doi.org/10.1145/3573105.3575671}{10.1145/3573105.3575671}.
Author manuscript:
\url{https://people.compute.dtu.dk/ahfrom/aesop-camera-ready.pdf}.

\bibitem{canonical}
Chase Norman and Jeremy Avigad.
\emph{Canonical for Automated Theorem Proving in Lean}.
2025. arXiv:2504.06239.
\url{https://arxiv.org/abs/2504.06239}.

\bibitem{canonicalmin}
Chase Norman and Jeremy Avigad.
\emph{Implementing Dependent Type Theory Inhabitation and Unification}.
March 2, 2026. arXiv:2603.01463.
\url{https://arxiv.org/html/2603.01463v1}.

\bibitem{probe}
Shraddha Barke, Hila Peleg, and Nadia Polikarpova.
Just-in-Time Learning for Bottom-Up Enumerative Synthesis.
\emph{Proceedings of the ACM on Programming Languages}, 4(OOPSLA), 2020.
DOI: \href{https://doi.org/10.1145/3428295}{10.1145/3428295}.
\url{https://arxiv.org/abs/2010.08663}.

\bibitem{luby}
Michael Luby, Alistair Sinclair, and David Zuckerman.
Optimal Speedup of Las Vegas Algorithms.
\emph{Information Processing Letters}, 47:173--180, 1993.
DOI: \href{https://doi.org/10.1016/0020-0190(93)90029-9}{10.1016/0020-0190(93)90029-9}.

\bibitem{whenfold}
Jeremy Gibbons, Graham Hutton, and Thorsten Altenkirch.
When is a Function a Fold or an Unfold?
\emph{Electronic Notes in Theoretical Computer Science}, 44(1), 2001.
Author publication record:
\url{https://www.cs.ox.ac.uk/publications/publication2368-abstract.html}.
Related author slides used for the explicit reverse example:
\url{https://people.cs.nott.ac.uk/pszgmh/wgp01/gibbons-slides.pdf}.

\bibitem{angluin}
Dana Angluin.
Learning Regular Sets from Queries and Counterexamples.
\emph{Information and Computation}, 75(2):87--106, 1987.
DOI: \href{https://doi.org/10.1016/0890-5401(87)90052-6}{10.1016/0890-5401(87)90052-6}.

\bibitem{wadler}
Philip Wadler.
Theorems for Free!
\emph{Functional Programming Languages and Computer Architecture}, 1989.
DOI: \href{https://doi.org/10.1145/99370.99404}{10.1145/99370.99404}.

\bibitem{autolifter}
Ruyi Ji, Yuwei Zhao, Yingfei Xiong, Di Wang, Lu Zhang, and Zhenjiang Hu.
\emph{Decomposition-Based Synthesis for Applying Divide-and-Conquer-Like
Algorithmic Paradigms}. arXiv:2202.12193, revision 4, January 29, 2024;
accepted for ACM TOPLAS as noted by that revision. Especially the algorithm,
minimality discussion, and limitations.
\url{https://arxiv.org/html/2202.12193v4}.

\bibitem{hazel}
Cyrus Omar, Ian Voysey, Ravi Chugh, and Matthew A. Hammer.
Live Functional Programming with Typed Holes.
\emph{Proceedings of the ACM on Programming Languages}, 3(POPL), 2019.
DOI: \href{https://doi.org/10.1145/3290327}{10.1145/3290327}.
\url{https://arxiv.org/html/1805.00155v4}.

\bibitem{smyth}
Justin Lubin, Nick Collins, Cyrus Omar, and Ravi Chugh.
Program Sketching with Live Bidirectional Evaluation.
\emph{Proceedings of the ACM on Programming Languages}, 4(ICFP), article 109,
2020. DOI: \href{https://doi.org/10.1145/3408991}{10.1145/3408991}.
The implementation discussion includes higher-order examples and polymorphism.
\url{https://arxiv.org/abs/1911.00583}.

\bibitem{burst}
Anders Miltner, Adrian Trejo Nu\~nez, Ana Brendel, Swarat Chaudhuri, and Isil Dillig.
Bottom-up Synthesis of Recursive Functional Programs using Angelic Execution.
\emph{Proceedings of the ACM on Programming Languages}, 6(POPL), 2022.
DOI: \href{https://doi.org/10.1145/3498682}{10.1145/3498682}.
\url{https://arxiv.org/abs/2107.06253}.

\bibitem{lean4lean}
Mario Carneiro.
\emph{Lean4Lean: Verifying a Typechecker for Lean, in Lean}.
arXiv:2403.14064, revision 3, September 14, 2025.
\url{https://arxiv.org/abs/2403.14064}.

\bibitem{leanmanual}
Lean development team.
\emph{The Lean Language Reference}, version 4.34.0. Especially the tactic
reference entries for induction, functional induction, simplification,
arithmetic, and general proof automation.
\url{https://lean-lang.org/doc/reference/4.34.0/Tactic-Proofs/Tactic-Reference/}.

\bibitem{leaninduction}
Lean development team.
\emph{Induction metaprogramming implementation}, version 4.34.0.
\path{src/Lean/Meta/Tactic/Induction.lean}. Includes
\code{getMajorTypeIndices}, \code{mkRecursorAppPrefix}, and
\code{Lean.MVarId.induction}.
\url{https://github.com/leanprover/lean4/blob/v4.34.0/src/Lean/Meta/Tactic/Induction.lean}.

\bibitem{synquid}
Nadia Polikarpova, Ivan Kuraj, and Armando Solar-Lezama.
Program Synthesis from Polymorphic Refinement Types.
\emph{PLDI}, 2016.
DOI: \href{https://doi.org/10.1145/2908080.2908093}{10.1145/2908080.2908093}.
\url{https://arxiv.org/abs/1510.08419}.

\bibitem{tygar}
Zheng Guo, Michael James, David Justo, Jiaxiao Zhou, Ziteng Wang,
Ranjit Jhala, and Nadia Polikarpova.
Program Synthesis by Type-Guided Abstraction Refinement.
\emph{Proceedings of the ACM on Programming Languages}, 4(POPL), article 12,
2020. DOI: \href{https://doi.org/10.1145/3371080}{10.1145/3371080}.
Especially Section 5.1 on type classes and implementation limitations.
\url{https://ranjitjhala.github.io/static/hoogle_plus.pdf}.

\bibitem{leansuggestions}
Lean development team.
\emph{An API for library suggestion algorithms}, version 4.34.0.
\path{src/Lean/LibrarySuggestions/Basic.lean}. Includes
\code{Selector}, \code{Config.hint}, and the documented absence of a default
library suggestion engine.
\url{https://github.com/leanprover/lean4/blob/v4.34.0/src/Lean/LibrarySuggestions/Basic.lean}.

\bibitem{leanhammer}
Thomas Zhu, Joshua Clune, Jeremy Avigad, Albert Q. Jiang, and Sean Welleck.
\emph{Premise Selection for a Lean Hammer}.
ICLR 2026; arXiv:2506.07477, revision 2, February 25, 2026.
\url{https://arxiv.org/html/2506.07477v2}.

\end{thebibliography}
\end{document}
